% !TeX root = main.tex
%% ===========================================================================
%%  config.tex -- YOUR DATA GOES HERE
%% ---------------------------------------------------------------------------
%%  Fill in every field below. Leave a field empty, e.g. \ThesisSubtitle{},
%%  and it simply disappears from the cover and the other pages.
%%  Use the official names of your university, school and degree exactly as
%%  they appear in your university's regulations (an official English name
%%  is usually published on the degree's web page).
%% ===========================================================================

% Optional: start from the style of a specific university (see presets/).
% Your own settings below take precedence over the preset.
% % !TeX root = ../main.tex
%% ===========================================================================
%%  Preset: Universidad de Jaén -- Escuela Politécnica Superior de Jaén (EPSJ)
%% ---------------------------------------------------------------------------
%%  Reproduces the style rules of the original template by Francisco Charte
%%  Ojeda (https://github.com/fcharte/TFGLatexTemplate), written for the EPSJ
%%  style guidelines: A4, 2.5 cm margins, 1.5 line spacing, first-line
%%  indentation, Arial 12 pt, header with the student's name and the thesis
%%  title, footer with the school's name and the page number.
%%
%%  Those guidelines were written for Bachelor's theses (TFG): check the
%%  specific rules of your Master's programme before relying on them.
%%
%%  How to use it:
%%    1. In main.tex, use the class options:  font=sans, uppercase
%%    2. At the top of config.tex add:        % !TeX root = ../main.tex
%% ===========================================================================
%%  Preset: Universidad de Jaén -- Escuela Politécnica Superior de Jaén (EPSJ)
%% ---------------------------------------------------------------------------
%%  Reproduces the style rules of the original template by Francisco Charte
%%  Ojeda (https://github.com/fcharte/TFGLatexTemplate), written for the EPSJ
%%  style guidelines: A4, 2.5 cm margins, 1.5 line spacing, first-line
%%  indentation, Arial 12 pt, header with the student's name and the thesis
%%  title, footer with the school's name and the page number.
%%
%%  Those guidelines were written for Bachelor's theses (TFG): check the
%%  specific rules of your Master's programme before relying on them.
%%
%%  How to use it:
%%    1. In main.tex, use the class options:  font=sans, uppercase
%%    2. At the top of config.tex add:        % !TeX root = ../main.tex
%% ===========================================================================
%%  Preset: Universidad de Jaén -- Escuela Politécnica Superior de Jaén (EPSJ)
%% ---------------------------------------------------------------------------
%%  Reproduces the style rules of the original template by Francisco Charte
%%  Ojeda (https://github.com/fcharte/TFGLatexTemplate), written for the EPSJ
%%  style guidelines: A4, 2.5 cm margins, 1.5 line spacing, first-line
%%  indentation, Arial 12 pt, header with the student's name and the thesis
%%  title, footer with the school's name and the page number.
%%
%%  Those guidelines were written for Bachelor's theses (TFG): check the
%%  specific rules of your Master's programme before relying on them.
%%
%%  How to use it:
%%    1. In main.tex, use the class options:  font=sans, uppercase
%%    2. At the top of config.tex add:        \input{presets/uja-epsj}
%% ===========================================================================
\University{Universidad de Jaén}
\Faculty{Escuela Politécnica Superior de Jaén}
\City{Jaén}
\FooterText{Escuela Politécnica Superior de Jaén}
% Orange of the EPSJ cover used in the original template
\definecolor{tfmaccent}{RGB}{198,97,21}

% Header: student's name and short title. Footer: school and page number.
\fancyhf{}
\fancyhead[LE,LO]{\footnotesize\color{tfmgray}\thesisauthor}
\fancyhead[RE,RO]{\footnotesize\color{tfmgray}\thesisshorttitle}
\fancyfoot[LE,LO]{\footnotesize\color{tfmgray}Escuela Politécnica Superior de Jaén}
\fancyfoot[RE,RO]{\small\thepage}
\renewcommand{\footrulewidth}{1pt}

%% ===========================================================================
\University{Universidad de Jaén}
\Faculty{Escuela Politécnica Superior de Jaén}
\City{Jaén}
\FooterText{Escuela Politécnica Superior de Jaén}
% Orange of the EPSJ cover used in the original template
\definecolor{tfmaccent}{RGB}{198,97,21}

% Header: student's name and short title. Footer: school and page number.
\fancyhf{}
\fancyhead[LE,LO]{\footnotesize\color{tfmgray}\thesisauthor}
\fancyhead[RE,RO]{\footnotesize\color{tfmgray}\thesisshorttitle}
\fancyfoot[LE,LO]{\footnotesize\color{tfmgray}Escuela Politécnica Superior de Jaén}
\fancyfoot[RE,RO]{\small\thepage}
\renewcommand{\footrulewidth}{1pt}

%% ===========================================================================
\University{Universidad de Jaén}
\Faculty{Escuela Politécnica Superior de Jaén}
\City{Jaén}
\FooterText{Escuela Politécnica Superior de Jaén}
% Orange of the EPSJ cover used in the original template
\definecolor{tfmaccent}{RGB}{198,97,21}

% Header: student's name and short title. Footer: school and page number.
\fancyhf{}
\fancyhead[LE,LO]{\footnotesize\color{tfmgray}\thesisauthor}
\fancyhead[RE,RO]{\footnotesize\color{tfmgray}\thesisshorttitle}
\fancyfoot[LE,LO]{\footnotesize\color{tfmgray}Escuela Politécnica Superior de Jaén}
\fancyfoot[RE,RO]{\small\thepage}
\renewcommand{\footrulewidth}{1pt}


%% --- The thesis --------------------------------------------------------------
% Full title: it must match the title of the approved thesis proposal
\ThesisTitle{Title of Your Master's Thesis, Exactly as in the Approved Proposal}
% Optional subtitle
\ThesisSubtitle{An optional subtitle that adds context}
% Short title for the page headers (defaults to the full title if empty)
\ThesisShortTitle{Short title for the page headers}
% Title in the second language (Spanish, Catalan, Galician, Basque...).
% Most universities require it when the thesis is written in English.
\ThesisTitleSecondLanguage{[Title of the thesis in the second language]}
% Kind of work. Change it if your university uses another English name
\ThesisType{Master's Thesis}
% Name of the work in the second language. The default depends on the
% secondlanguage class option (e.g. "Trabajo Fin de Máster" for spanish).
% Uncomment to override it, or leave it empty to hide it on the cover:
% \ThesisTypeSecondLanguage{Trabajo Fin de Máster}

%% --- The author -------------------------------------------------------------
\ThesisAuthor{First-name Surname1 Surname2}
\AuthorEmail{student@university.es}      % optional, shown on the licence page
\AuthorIDNumber{}                        % optional (DNI/NIE) for the declarations

%% --- Supervisors ------------------------------------------------------------
% \Supervisor[Role shown on the cover]{Full name}{Affiliation}
% Repeat the command for every supervisor (co-supervisors, external
% supervisors from a company, etc.). The role defaults to "Supervisor".
\Supervisor{Dr. First-name Surname}{Department of Computer Science, University Name}
\Supervisor[Co-supervisor]{Dr. First-name Surname}{Research Institute or Company}

%% --- Institution and degree -------------------------------------------------
\University{Universidad de [University Name]}
\Faculty{[School or Faculty Name]}
\Department{[Department of the supervisor]}
\Degree{Master's Degree in [Name of the Master's Programme]}
\City{[City]}
\AcademicYear{2025/2026}
% Submission date as it should be printed. \today prints the compilation date.
\SubmissionDate{June 2026}
% Optional: defence session ("Ordinary session", "July 2026", ...)
\DefenceSession{}

%% --- Keywords ---------------------------------------------------------------
% Usually 3 to 6, separated by semicolons
\Keywords{Master's thesis; LaTeX; template; Spanish universities; reproducibility}
\KeywordsSecondLanguage{[keyword 1; keyword 2; keyword 3]}

%% --- Logos and cover --------------------------------------------------------
% Put your logos in figures/logos/. While the file does not exist, a neutral
% placeholder is drawn instead. Vector formats (PDF, EPS) look best.
% Download the logo from your university's corporate identity web page.
\UniversityLogo{figures/logos/university-logo.pdf}
\SecondaryLogo{}                     % e.g. school/faculty logo (optional)
% If your university provides an OFFICIAL COVER as a PDF (filled in with your
% data), point to it here: it becomes page 1 and the generated cover follows
% as the title page. Leave it empty to use only the generated cover.
\OfficialCover{}

%% --- Colours -----------------------------------------------------------------
% Accent colour of the cover, headings, captions and links. Use the corporate
% colour of your university if you like (any HTML hex code works).
\definecolor{tfmaccent}{HTML}{1F4E79}

%% --- Licence -----------------------------------------------------------------
% all-rights-reserved | by | by-sa | by-nc | by-nc-sa | by-nd | by-nc-nd
% Institutional repositories in Spain often suggest by-nc-nd. Check yours.
\ThesisLicence{by-nc-nd}

%% --- Optional pages ----------------------------------------------------------
% true = include the page, false = leave it out
\settoggle{tfmLicencePage}{true}      % copyright / Creative Commons notice
\settoggle{tfmApprovalPage}{true}     % supervisors' approval for the defence
\settoggle{tfmOriginalityPage}{true}  % declaration of originality / authorship
\settoggle{tfmTemplateGuide}{true}    % the "Using This Template" chapter.
                                      % Set it to false when you start writing!

%% --- Footer ------------------------------------------------------------------
% Text in the footer of every page (defaults to the university name)
% \FooterText{Faculty of Sciences -- University Name}

%% --- Layout overrides (only if your regulations demand something else) --------
% The defaults are: A4, 2.5 cm margins (3 cm inner margin in print mode),
% 1.5 line spacing, 12 pt text, first-line indentation.
% \geometry{top=3cm, bottom=3cm, left=3cm, right=2.5cm}
% \setstretch{1.25}
% \setlength{\parindent}{0pt}
% \setlength{\parskip}{0.8em}
% Headers and footers can be redefined with fancyhdr, e.g.:
% \fancyhead[L]{\footnotesize\color{tfmgray}\thesisauthor}
